%%%%%%%%%%%%%%%%%%%%%%%%%%%%%%%%%%%%%%%%%%%%%%%%%%%%%%
\chapter{Linear Algebra in Computer Processor Design}
%%%%%%%%%%%%%%%%%%%%%%%%%%%%%%%%%%%%%%%%%%%%%%%%%%%%%%

%*********************
\section{Introduction}
%*********************
Linear algebra is a foundational element of modern computing, underpinning everything from scientific simulations to machine learning algorithms. At its core, linear algebra deals with vectors (one-dimensional arrays of numbers), matrices (two-dimensional arrays), and higher-dimensional analogues often called \textit{tensors}. In computing terms, a \textbf{vector} can be a simple list of data (e.g. a sequence of signal samples or a coordinate vector), a \textbf{matrix} a table of data (e.g. a 2D image or a transformation matrix), and a \textbf{tensor} a multi-dimensional array (common in machine learning to represent batches of multi-dimensional data). Processor architectures have evolved to handle these structures efficiently because so many computational tasks boil down to operations on vectors, matrices, and tensors. In fact, operations like the dot product (for vectors), matrix multiplication, and even more complex tensor contractions are central to algorithms in digital signal processing, computer graphics, and artificial intelligence.

This chapter explores how linear algebra concepts have influenced and shaped computer processor design. We will see that different processor types have been optimized for different dimensions of linear algebra: traditional \textbf{Digital Signal Processors (DSPs)} excel at vector (1D) operations, modern \textbf{Graphics Processing Units (GPUs)} are tailored for matrix (2D) operations, and recent specialized accelerators like \textbf{Tensor Processing Units (TPUs)} target high-dimensional tensor operations (3D and beyond). We will trace the historical development of these ideas, from early supercomputers that introduced vector processing, through the advent of GPUs for graphics (and later general-purpose computing), to the current landscape of AI accelerators. We will also examine how contemporary processor instruction sets (with an emphasis on RISC-V, ARM, and x86 architectures) incorporate linear algebra operations, and how typical linear algebra computations (dot products, matrix multiplies, tensor contractions, etc.) map onto hardware structures and instructions. Furthermore, we discuss how higher-rank tensors can often be transformed or decomposed into lower-rank forms to make computations more efficient, and how such techniques are leveraged in hardware and software optimizations. Finally, we consider software-hardware co-design aspects – how software libraries, compilers, and frameworks collaborate with hardware features to achieve high performance on linear algebra tasks – and speculate on future directions in processor architecture (such as analog, photonic, and quantum computing) with respect to linear algebra.

%****************************************************************
\section{Historical Development of Linear Algebra in Processors} 
%****************************************************************
Linear algebra computations have been important in computing since the early days of digital computers. Early scientific computing in the 1960s and 1970s already relied on solving linear equations, performing matrix factorizations, etc., which are heavy uses of linear algebra. This drove the development of optimized software libraries (like the precursor of today’s BLAS and LINPACK) and also influenced hardware design. A landmark in this history was the introduction of \textbf{vector processors} in supercomputers. In 1976, Seymour Cray’s \emph{Cray-1} supercomputer became the first commercially successful machine with vector processing capabilities【42†L19-L27】. The Cray-1 could load sets of data into special \emph{vector registers} and operate on these long vectors with a single instruction, dramatically accelerating linear algebra operations common in science and engineering. In the Cray-1, each vector register held 64 elements of 64-bit data, and a dedicated \emph{vector length} register could be set to operate on varying vector sizes【26†L1-L8】【41†L1-L4】. By keeping vectors in fast registers and using pipelined arithmetic units, the Cray-1 achieved about 80 MFLOPS (million floating-point operations per second) sustained and up to 240 MFLOPS peak【26†L5-L13】—unprecedented performance at the time, largely thanks to efficient vectorized linear algebra. This established the paradigm that to go faster, computers needed to process multiple data elements in parallel using a single instruction stream, a technique directly beneficial for linear algebra. Subsequent supercomputers (from CDC, NEC, etc.) and array processors in the 1980s continued this trend of vectorization.

While supercomputers led the way, microprocessor designs also started to adopt linear algebra oriented features as those workloads became more common in broader applications (such as signal processing and multimedia). \textbf{Digital Signal Processors (DSPs)} emerged in the late 1970s and 1980s as specialized microprocessors for real-time signal processing (audio, communications, etc.), and they heavily relied on the linear algebra operation of the dot product (multiply-accumulate). DSP architectures introduced hardware MAC (Multiply-Accumulate) units that could multiply two numbers and add the result to an accumulator in one step, effectively computing a dot-product term in each cycle【8†L133-L142】【8†L149-L157】. In fact, the first modern processors to include dedicated MAC units were DSPs【8†L149-L157】. This capability was soon incorporated into general-purpose processors as well, since many algorithms beyond signal processing benefit from fast multiply-add operations. By the 1990s and early 2000s, mainstream CPU architectures (x86, ARM, etc.) introduced their own single-instruction multiple-data (\textbf{SIMD}) extensions – such as Intel’s MMX and SSE, and ARM’s NEON – allowing one instruction to perform arithmetic on several vector elements in parallel. These SIMD extensions were initially modest (processing, say, 4 elements at a time), but they marked the spread of vector processing concepts into commodity computing. An example is Intel’s SSE (Streaming SIMD Extensions) in 1999, which could perform four 32-bit floating-point operations in parallel within one 128-bit register. This trend continued with wider vectors (Intel’s AVX later widened this to 256 bits and AVX-512 to 512 bits, as we will discuss). The motivation for these was partly graphics and multimedia (e.g. encoding/decoding, where linear transformations are common) and partly scientific computing – all essentially linear algebra computations on arrays of data.

In parallel, the 1990s saw the rise of \textbf{Graphics Processing Units (GPUs)}. GPUs were originally designed to accelerate the linear algebra of graphics: operations like transforming 3D coordinates (which use $4\times4$ matrices), lighting calculations (dot products between vectors), and texture mappings. Early GPUs had fixed-function pipelines for these linear algebra operations, but over time they evolved into programmable processors with many cores capable of general numerical computation. By the mid-2000s, GPUs became general-purpose enough (through shader programming and later CUDA/OpenCL frameworks) that researchers started using them for non-graphics tasks, especially those involving linear algebra on large matrices. In 2007 NVIDIA introduced \emph{CUDA}, which allowed developers to program GPUs for arbitrary data-parallel tasks. One of the first and most impactful use-cases was accelerating dense matrix computations and linear algebra routines (like matrix multiplication, solvers, FFTs) for scientific computing. GPUs were attractive because, compared to CPUs, they devoted a much larger proportion of transistors to arithmetic units (ALUs) rather than caches or control, yielding much higher raw FLOPS (floating point operations per second) on matrix calculations【46†L7-L15】. A modern GPU consists of an array of hundreds or thousands of simpler cores that can execute threads in parallel, using a Single-Instruction Multiple-Thread (SIMT) model to apply the same instructions across many data elements. This essentially makes a GPU a massive vector processor: indeed, modern GPUs “include an array of shader pipelines…and can be considered vector processors”【5†L25-L33】, employing similar strategies as classic vector machines to hide memory latency (through thread switching rather than deep pipelines). This shift was significant in the historical trajectory: it brought supercomputer-level linear algebra acceleration (once only in Cray-type machines) to commodity hardware. By 2010s, GPUs became the workhorse for machine learning training exactly because neural network training involves tons of linear algebra (matrix multiplications for fully-connected layers, tensor convolutions, etc.), and GPUs could perform these in parallel at high speed.

The explosion of \textbf{machine learning (ML)} and especially deep learning in the 2010s put linear algebra at the center of computing arguably more than ever before. Neural networks are essentially built on linear algebra operations – for example, each layer of a neural network can be viewed as multiplying an input vector by a weight matrix, accumulating results, and applying simple nonlinear functions. Convolutional layers in vision can be transformed into matrix-multiplication as well, or viewed as tensor operations. As models grew, so did the demand for accelerating these operations. This led to a new class of processors and accelerators in the late 2010s: \textbf{Tensor Processing Units (TPUs)} and various Neural Processing Units (NPUs). Google’s TPU (first announced in 2016) is a seminal example, designed specifically to accelerate the kinds of matrix and vector operations used in ML. In essence, Google “reincarnated the systolic array as the TPU” to meet the needs of low-precision (e.g. 8-bit or 16-bit) multiplication of large matrices in neural network training and inference【32†L227-L234】. The first-generation TPU focused on inference and featured a large 256$\times$256 matrix multiply unit implemented as a systolic array, achieving extremely high throughput on matrix multiply operations while sacrificing general-purpose flexibility【32†L229-L237】. This philosophy – specialize hardware to do (almost) nothing but linear algebra, but do it blindingly fast – paid off for deep learning. Subsequent TPUs (v2, v3, etc.) expanded to handle training with floating-point (bfloat16) and scaled out to “pods” of many TPU chips. Other companies and research groups followed suit, producing AI accelerators that are essentially dense linear algebra engines (e.g. systolic arrays, massive MAC arrays, etc.). As a result, as of the mid-2020s, linear algebra has influenced processor design at all scales: from special functional units in general-purpose CPUs, to GPUs in every PC or phone, to cloud-scale accelerators for AI. The next sections discuss in detail vectors, matrices, and tensors in processor design, and how current hardware (especially RISC-V, ARM, and x86 architectures) implements and optimizes these operations.

%************************************************
\section{Mathematical Foundations for Processors}
%************************************************
Although linear algebra is a vast subject, modern processors rely heavily on a specific subset of concepts: how vectors are represented and combined, and how to multiply vectors \emph{internally} or form products that produce higher-dimensional results. This section provides a concise review of those relevant mathematical tools and offers geometric insight into why they matter in hardware designs.

%---------------------------------------
\subsection{Vectors and Vector Addition}
%---------------------------------------
A \emph{vector} is typically represented in mathematics as an ordered list of numbers:
\[
\mathbf{x} = \begin{bmatrix} x_1 \\ x_2 \\ \vdots \\ x_n \end{bmatrix},
\]
where each $x_i$ is a component of the vector. We usually think of such vectors as points or directions in an $n$-dimensional space. In a geometric sense, a vector can be visualized as an arrow from the origin $(0,0,\ldots,0)$ to the point $(x_1, x_2, \ldots, x_n)$.

\paragraph{Vector Addition.} 
Adding two vectors of the same dimension,
\[
\mathbf{x} = \begin{bmatrix} x_1 \\ x_2 \\ \vdots \\ x_n \end{bmatrix}, 
\quad
\mathbf{y} = \begin{bmatrix} y_1 \\ y_2 \\ \vdots \\ y_n \end{bmatrix},
\]
yields another $n$-dimensional vector:
\[
\mathbf{x} + \mathbf{y} 
= \begin{bmatrix} x_1 + y_1 \\ x_2 + y_2 \\ \vdots \\ x_n + y_n \end{bmatrix}.
\]
Geometrically, this is the “tip-to-tail” rule: place $\mathbf{y}$ so that its tail meets the tip of $\mathbf{x}$. The resulting arrow from the origin to the final tip is $\mathbf{x} + \mathbf{y}$. 

\paragraph{Physical Interpretation.} 
In processor hardware, vector addition corresponds to adding element-wise data in parallel. If the processor has \emph{SIMD} instructions (e.g., AVX, NEON), it can add several pairs of elements at once. On a Digital Signal Processor (DSP), vector addition might be pipelined or combined with multiply-accumulate operations for filtering tasks. Although the geometric interpretation is not typically “visualized” in hardware, the principle of adding corresponding components exactly translates to how the hardware sums each register lane.

%--------------------------------
\subsection{Inner (Dot) Product}
%--------------------------------
The \textbf{inner product} (often called the dot product in $\mathbb{R}^n$) of two vectors of the same length $n$ is
\[
\mathbf{x} \cdot \mathbf{y} \;=\; \sum_{i=1}^{n} x_i \, y_i.
\]
Geometrically, this measures how much $\mathbf{x}$ and $\mathbf{y}$ align. In particular,
\[
\mathbf{x} \cdot \mathbf{y} \;=\; \|\mathbf{x}\|\|\mathbf{y}\|\cos\theta,
\]
where $\theta$ is the angle between the two vectors, and $\|\mathbf{x}\|$ denotes the length (Euclidean norm) of $\mathbf{x}$.

\paragraph{Processor Relevance.} 
In physical terms, the dot product is the sum of pairwise multiplications. Digital Signal Processors exploit a \textbf{multiply-accumulate (MAC)} unit to execute one term of the dot product in each cycle, effectively computing $x_i \times y_i + \text{accumulator}$. Modern CPUs provide \emph{fused multiply-add (FMA)} instructions that bundle the multiplication and addition into a single instruction. In vectorized form (SIMD), the hardware can multiply several pairs in parallel and then \emph{horizontally} add them to get a final scalar. 

\paragraph{Geometric Benefit.} 
The hardware’s job is to replicate this fundamental operation many times, as so many algorithms rely on correlations or projections (both concepts intimately connected to dot products). Think of a neural network layer: the \texttt{dot} between a weight vector and an input vector is the main arithmetic in computing the next layer’s activations. GPUs and AI accelerators further unify these operations to handle millions (or billions) of inner products every second by distributing them across their parallel cores.

%--------------------------
\subsection{Outer Product}
%-------------------------
Whereas the dot product of two vectors yields a scalar, the \textbf{outer product} of a column vector $\mathbf{x}\in \mathbb{R}^m$ with a row vector $\mathbf{y}^T\in \mathbb{R}^n$ produces an $m\times n$ matrix:
\[
\mathbf{x}\,\mathbf{y}^T 
\;=\; \begin{bmatrix}
x_1\,y_1 & x_1\,y_2 & \ldots & x_1\,y_n \\[6pt]
x_2\,y_1 & x_2\,y_2 & \ldots & x_2\,y_n \\[3pt]
\vdots   & \vdots   & \ddots & \vdots   \\[3pt]
x_m\,y_1 & x_m\,y_2 & \ldots & x_m\,y_n
\end{bmatrix}.
\]
Notice each entry is simply $x_i \cdot y_j$.

\paragraph{Geometric and Physical Meaning.} 
Geometrically, the outer product represents a rank-1 map from $\mathbb{R}^n$ to $\mathbb{R}^m$. It’s “rank-1” because any column of this matrix is just a scaled version of $\mathbf{x}$, and any row is a scaled version of $\mathbf{y}^T$. In hardware and numerical linear algebra, outer products appear in algorithms that update matrices with a “rank-1 update.” For instance, in block-based matrix multiply, if you take one row block of $A$ (seen as a column vector) and one column block of $B$ (as a row vector), their outer product can be added into the sub-block of $C$. 

\paragraph{Why Processors Care.} 
Many matrix-multiply kernels use an \emph{outer-product formulation} to compute $C += A \times B$ by repeatedly forming small outer products. Each outer product is highly parallelizable: the hardware just needs to multiply each pair $(x_i, y_j)$ and store or accumulate the result in the correct position. Specialized instructions (like Intel’s AMX tile or ARM’s SME) effectively perform blocks of outer products in hardware. This approach also relates to \emph{tensor} operations, where repeated outer products can represent low-rank approximations or partial expansions.

%-------------------
\subsection{Tensors}
%-------------------
A \emph{tensor} is a mathematical structure that generalizes vectors and matrices to higher dimensions. 
In the same way that a vector can be viewed as a one-dimensional array of elements, and a matrix as 
a two-dimensional arrangement of elements, an $n$-th order tensor is an $n$-dimensional array 
of numbers. Concretely, a rank-3 tensor might be written as
\[
\mathcal{T}_{i,j,k}, \quad i \in \{1,\dots,I\}, \; j \in \{1,\dots,J\}, \; k \in \{1,\dots,K\},
\]
which one can envision as a ``stack'' of $I \times J$ matrices indexed by $k$. Higher-rank tensors 
(i.e., more than three dimensions) are typically seen in machine learning applications, where one 
dimension may correspond to the \emph{batch} of input data, another to the \emph{height} of an image, 
another to the \emph{width}, another to \emph{channels}, and so forth.

\paragraph{Geometric and Physical Interpretation.}
From a geometric standpoint, a rank-1 tensor in $\mathbb{R}^n$ is simply a vector, 
while a rank-2 tensor (in the linear-algebra sense) is just a matrix. Once we go to rank 3 or higher, 
it becomes more complicated to visualize, but the essence remains: each index corresponds to a distinct 
``direction'' or mode in the data. In physics, for example, a second-rank tensor might describe stress 
or strain in materials, mapping one vector (a direction of force) to another (the resulting deformation). 
In image processing and deep learning, a 4D tensor can represent multiple images, each with a particular 
height, width, and color-channel dimension.

\paragraph{Basic Tensor Operations.}
Many tensor operations boil down to repeated \emph{vector} or \emph{matrix} operations. 
For instance, a \emph{tensor contraction} over one or more indices is a higher-dimensional analogue 
of the dot product, summing products of elements that share a common index. A 3D tensor $\mathcal{T}_{i,j,k}$ 
contracted with another 3D tensor $\mathcal{U}_{j,k,\ell}$ over indices $j$ and $k$ might yield a 
2D matrix $M_{i,\ell}$:
\[
M_{i,\ell} \;=\;\sum_{j,k}\; \mathcal{T}_{i,j,k}\,\mathcal{U}_{j,k,\ell}.
\]
Such a contraction is analogous to nested loops of multiply-accumulate operations on the hardware side. 

\paragraph{Processor-Level Relevance.}
In computer processors, especially in AI accelerators, tensors are stored as multi-dimensional 
arrays in memory, often laid out contiguously in row-major or other specialized formats. 
Hardware typically decomposes or ``flattens'' higher-rank tensor operations into sequences of 
\emph{matrix} or \emph{vector} operations, because core instructions for multiplication 
and addition are designed around 1D or 2D data. As an example, convolutional layers in neural 
networks (which rely on 4D or 5D tensors of data and kernel weights) are usually transformed 
into large matrix-matrix multiplications internally (via \textit{im2col} or similar reshaping). 
Hence, although tensors are inherently high-dimensional, their computations are translated into 
the 1D or 2D operations that processors (CPU, GPU, or specialized tensor hardware) already optimize.

\paragraph{Physical Insight.}
One way to appreciate higher-rank tensors is to see them as ``multi-linear maps.'' 
A vector in $\mathbb{R}^n$ can be seen as a linear map from $\mathbb{R}^1$ to $\mathbb{R}^n$, 
and a matrix (as a second-rank tensor) is a linear map from $\mathbb{R}^n$ to $\mathbb{R}^m$. 
By extension, an $n$-th rank tensor is a map that can accept $n$ vectors, each possibly from a 
different space, and produce one output (scalar, vector, or another tensor). Physically, 
this underpins phenomena like elasticity (a rank-4 tensor describing stress-strain relationships), 
or multi-channel images in vision (4D or 5D structures). In digital hardware, while there is 
no direct geometric representation of these objects, the multi-index structure is preserved 
in \emph{address arithmetic}: the hardware or compiler determines how to iterate over all 
tensor indices and feed the appropriate multiply-add pipelines accordingly. 
This is why processor vendors introduce specialized \emph{tensor} instructions (like ``tensor 
cores'' or ``systolic arrays'') to handle multi-dimensional data more directly, reflecting the 
mathematical fact that so many workloads rely on these higher-order interactions.

%---------------------------------------------------------
\subsection{Physical Interpretation in Processors Summary}
%----------------------------------------------------------
From a hardware perspective, these mathematical constructs translate to:
\begin{itemize}
  \item \textbf{Vector addition:} Summation of array lanes within specialized registers in parallel.
  \item \textbf{Dot product:} Summation of pairwise products, often accelerated by fused multiply-add pipelines (MAC or FMA), possibly over wide SIMD lanes (e.g., 8, 16, or more multiplications at once).
  \item \textbf{Outer product:} Generating a matrix from two vectors, each element computed by a multiply. In hardware, one can think of a set of multipliers that take each element of $\mathbf{x}$ and multiply by each element of $\mathbf{y}$ in a systematic pattern, storing or accumulating these partial products in a 2D array of registers or memory.
  \item \textbf{Tensor Contraction:} In higher-rank tensors, a \emph{contraction} is essentially a 
multi-dimensional generalization of the dot product, summing products of elements that share certain 
indices. At the hardware level, this breaks down into nested multiply-add operations over the 
contracted dimensions. Modern processors typically handle tensor contractions by reshaping them into 
matrix or vector operations (such as GEMM calls) that exercise the same fused multiply-add pipelines 
and caching strategies employed for dot products and outer products.```

\end{itemize}

These elemental operations (vector addition, dot products, outer products) recur in nearly every linear algebra routine---from basic matrix multiply to complex tensor factorizations. Hence, processor designers add ever more sophisticated instructions and specialized hardware blocks to implement them efficiently. At the algorithmic level, the geometry of “directions,” “magnitudes,” and “projections” becomes a set of parallel multiply-add loops at the silicon level. Despite the dry mechanical approach inside the hardware, the underlying mathematics preserves all the elegant geometric interpretations: angles, orthogonal projections, expansions by rank-1 updates, and so on.


%******************************************************
\section{Vectors and Digital Signal Processors (DSPs)} 
%******************************************************
A \textbf{vector} in computational terms is an ordered one-dimensional set of data. Common operations on vectors include element-wise arithmetic, scaling by a scalar, and the dot product (or inner product). The dot product of two equal-length vectors $x = (x_1,\dots,x_n)$ and $y = (y_1,\dots,y_n)$ is a single number computed as: 
\[ 
x \cdot y = \sum_{i=1}^{n} x_i \, y_i \,,
\] 
essentially a sum of pairwise multiplications. This operation is ubiquitous: it appears in signal processing (as a finite impulse response filter convolution, which is a sliding dot product), in graphics (to compute lighting or angles between vectors), in machine learning (to compute neuron activations), and many other domains. Because the dot product involves performing the same multiply-add operation on every element of a vector, it is a prime candidate for parallelization in hardware.

%---------------------------------------
\subsection{Vector Processing Concepts} 
%--------------------------------------
A processor that can handle vectors efficiently typically includes hardware to apply the same operation to multiple elements at once. In a classical DSP or vector processor, this is achieved with a combination of: (1) \emph{vector registers} that can hold an array of numbers, and (2) \emph{vector instructions} that operate on those entire registers. For example, a vector processor might have an instruction to multiply two length-64 vectors (from two vector registers) element-wise, producing 64 results in a vector register, and another instruction to sum (reduce) a vector register’s contents to a scalar (accumulate). Traditional DSP chips often took a slightly different but related approach: because their primary workload was dot products, they included a \textbf{Multiply-Accumulate (MAC)} unit that could multiply two values and immediately add the result to an accumulator register in one instruction cycle. The MAC operation effectively does one iteration of a dot-product loop in hardware. For instance, if we wanted to compute a dot product $s = \sum_{i=1}^{n} a_i b_i$, a MAC-based loop would do something like:
\[
s := 0; \quad \text{for } i=1 \text{ to } n: \; s := s + (a_i \times b_i)~,
\] 
where each step uses the MAC unit. A single MAC instruction computes $a_i\times b_i + s_{\text{prev}}$ to produce $s_{\text{new}}$. 

Early DSPs were built around fast MAC units. As noted earlier, DSPs were the first class of processors to integrate dedicated hardware for multiply-accumulate【8†L149-L157】. For example, a classic DSP might have been able to start a new MAC operation every clock cycle, thus effectively computing one dot-product term per cycle (possibly with pipelining). Many DSPs also supported *dual* memory accesses per cycle (fetching both an $a_i$ and $b_i$ simultaneously) and had specialized addressing modes (to iterate through arrays or implement circular buffers for convolution buffers). All these features were aimed at pumping the MAC unit with data continuously, which is exactly what’s needed for vector arithmetic like dot products. A concrete example is the Texas Instruments TMS320 series (one of the earliest DSP families in the 1980s): it had an architecture optimized for repeated multiply-adds with minimal overhead, and could do a 16-bit multiply and 32-bit accumulate in a single cycle, with specialized hardware for shifting (for scaling) and saturating arithmetic (useful in fixed-point signal processing). In the context of linear algebra, such a DSP could be seen as a vector processor for 1-D vectors: it excelled at vector-vector operations.

General-purpose CPUs eventually adopted the MAC concept in the form of \textbf{fused multiply-add (FMA)} instructions. An FMA performs $d := a \times b + c$ in one go, which improves both speed and precision (since it does one rounding at the end for floating-point). Today, virtually all CPUs (x86, ARM, RISC-V, etc. with floating-point units) have FMA in hardware. This is a direct legacy of the influence of linear algebra: FMAs significantly speed up matrix and vector computations (a dot product of length $n$ can be done with $n$ FMA operations). As a historical note, the use of MAC/FMA in general computing was driven by workloads like digital media processing and scientific computing, and the IEEE-754 standard eventually added FMA as an optional operation due to its prevalence.

\subsection{Example: Vector Instruction Set Usage}  
Modern SIMD instruction sets (which are part of general-purpose ISAs like ARM, x86, RISC-V) allow a form of vector processing akin to a DSP. For instance, ARM’s NEON (available in virtually all smartphones) can handle vectors of 4 floats or 8 half-floats at a time. x86 has AVX which can handle vectors of 8 floats (in AVX2, 256-bit) or even 16 floats (in AVX-512, 512-bit) at once. To illustrate, consider adding two vectors of 8 single-precision numbers on an AVX2-capable CPU: instead of adding element-by-element in a loop, a single vector instruction can load 8 numbers into a 256-bit YMM register, another can load another 8 numbers into another register, and one \texttt{VADDPS} instruction (Vector Add Packed Single) will produce 8 results in one go. This yields a theoretical 8x speedup (minus some overhead) by utilizing data parallelism.

As a slightly more advanced example, ARMv8.2 introduced dedicated dot product instructions (UDOT/SDOT) which take four 8-bit elements from one vector and four 8-bit elements from another vector, multiply them pairwise and accumulate into a 32-bit result – effectively computing a 4-term dot product with one instruction【44†L11-L19】. Such instructions were added because of the importance of 8-bit integer arithmetic in ML (for neural network inference where weights and activations are often 8-bit). By doing 4 MACs at once, the dot-product throughput is multiplied. ARM’s Scalable Vector Extension (SVE) goes further by allowing vector lengths up to 2048 bits and also supports per-element predication (so it can handle vector lengths not a multiple of the hardware width gracefully). RISC-V’s Vector Extension similarly allows a single loop to operate on vectors of arbitrary length in chunks fitting the hardware, using a concept of \textit{vector length register} much like the Cray-1 had.

From a programming perspective, using these vector instructions might look like the following pseudocode (in RISC-V style) for a dot product of two 32-bit integer arrays, to illustrate how a compiler or assembly programmer might leverage them:

\begin{verbatim}
    asm
    vsetvli   t0, n, e32,m8    # Set vector length (VL) to handle 32-bit elems, 
                               # using 8 elements per vector operation (example)
    vle32.v   v1, (a0)         # Load vector v1 with elements from array A (pointer in a0)
    vle32.v   v2, (a1)         # Load vector v2 with elements from array B (pointer in a1)
    vmul.vv   v3, v1, v2       # Element-wise multiply v1 and v2 -> v3
    vredsum.vs v4, v3, v0      # Reduce (sum) vector v3 into scalar v4 (using v0 as initial)
    # at this point, v4[0] contains the partial sum of this segment of the dot product
\end{verbatim}

In a loop, the above would accumulate dot product segments. In reality, architectures like x86/ARM have fixed vector lengths (like 128-bit, 256-bit, etc.), whereas RISC-V and some ARM SVE implementations allow the hardware to choose a length (for instance, \texttt{m8} might mean 8 elements, but the actual hardware might use 128 or 256-bit registers – the compiler sets \texttt{m8} if it expects at most 8). The details are complex, but the takeaway is that vector instructions enable processors to handle 1D array operations (like our dot product) much faster by exploiting parallelism.

\subsection{DSPs Today and Integration in CPUs}  
While stand-alone DSP chips still exist for certain embedded applications, many of their features have been integrated into general CPUs or system-on-chip designs. It’s common to find a “DSP extension” in modern ISAs. For example, RISC-V has a proposed P-extension for DSP with features like 16-bit MACs, and ARM has the Helium extension (M-Profile Vector Extension) for microcontrollers which brings small-scale SIMD for DSP on microcontrollers. These are essentially vector (1D) accelerators. They accelerate not only classic signal processing but also general vector math (which is useful in linear algebra routines). In summary, vectors and their operations are the backbone of a vast range of computations, and processor designers have provided a wealth of features (from MAC units to wide SIMD registers) to make vector processing fast.

\section{Matrices and Graphics Processing Units (GPUs)}  
A \textbf{matrix} is a two-dimensional array of numbers, usually denoted with two indices (row and column). Matrices are central in linear algebra because they represent linear transformations and also because many higher-dimensional problems (like solving linear systems, least squares, eigenvalue problems, etc.) reduce to matrix computations. The single most important matrix operation is \textbf{matrix multiplication}. Multiplying an $m\times k$ matrix $A$ by a $k\times n$ matrix $B$ produces an $m\times n$ matrix $C$, where each element is a dot product of a row of $A$ with a column of $B$: 
\[ 
C_{i,j} = \sum_{\ell=1}^{k} A_{i,\ell} \times B_{\ell,j} \,.
\]
This operation involves $m \times n \times k$ multiply-accumulate operations, which for large matrices is extremely compute-intensive. Matrix multiplication (often abbreviated GEMM for “general matrix multiply”) forms the core of many algorithms in science and engineering; for instance, it is the key computation in solving systems of equations, in many signal processing transforms, and in virtually all deep neural networks (where matrix multiplies occur in fully-connected layers or as parts of convolution). It’s not an exaggeration to say that matrix multiplication is one of the most important workloads in computing. In fact, in modern AI workloads like large neural networks, over 90% of the compute is spent on matrix multiplications【37†L125-L134】.

\subsection{Why GPUs Excel at Matrix Operations}  
GPUs are designed to handle operations on large arrays of pixels, vertices, etc., which naturally translates to handling large matrices. A GPU contains hundreds or thousands of simple cores (ALUs) that can operate in parallel. While a CPU might have, say, 8 or 16 vector lanes that can perform 8 or 16 operations in parallel, a GPU can effectively execute thousands of operations in parallel by running many threads. For example, consider computing a matrix multiplication $C = A \times B$. This can be parallelized by computing each $C_{ij}$ (one dot product) on a separate thread, or by dividing the output matrix into tiles and assigning each tile to a thread block. Modern GPUs do exactly this: they “implement general matrix multiplication (GEMM) by partitioning the output matrix into tiles, which are then assigned to thread blocks”【11†L37-L44】. Each thread block will load portions of $A$ and $B$ into fast on-chip memory (shared memory), multiply-add them to accumulate a submatrix of $C$, and then write the result out to global memory. By tiling, GPUs ensure data reuse (each loaded chunk of $A$ and $B$ is used for many multiply-adds) and by using many thread blocks, they ensure all streaming multiprocessors (SMs) on the GPU are busy.

Another aspect that gives GPUs an edge is memory bandwidth and latency hiding. Matrix operations require a lot of data to be moved, and GPUs are built with very high bandwidth memory (e.g., GDDR or HBM memory, offering significantly higher throughput than typical CPU DRAM channels). They also hide latency by having many threads: if one group of threads is waiting for data from memory, another can use the ALUs, so the ALUs are rarely idle【11†L29-L37】【11†L49-L57】. This is the SIMT model: all threads in a “warp” execute the same instruction (like a vector add), but since each thread works on different data, it’s conceptually like operating on a matrix or large array in parallel. In effect, GPUs achieve massive parallel vector processing by using threads – often described as a difference between SIMD (single instruction, multiple data) vs SIMT (single instruction, multiple threads) architectures【5†L49-L57】【5†L27-L35】. In practice, both aim to do many operations with one instruction dispatch. A GPU warp of 32 threads might execute one vector instruction that covers 32 elements.

Historically, GPUs started to be used for matrix-heavy computations in the mid-2000s, and by the 2010s, libraries like cuBLAS (CUDA BLAS for NVIDIA GPUs) and cuDNN (for deep learning primitives) were highly optimized to make full use of GPU hardware for linear algebra. The speedups over CPUs were dramatic, often an order of magnitude or more, especially for large matrices where the GPU can amortize memory transfers and use its thousands of ALUs in parallel. This is why, for example, training a neural network on a GPU was found to be much faster than on CPU, enabling the recent breakthroughs in deep learning. It’s not that each GPU core is faster (in fact, each core is much slower than a CPU core), but that thousands of them working together on matrix operations achieve throughput a CPU can’t match【46†L1-L9】.

\subsection{GPU Hardware Features for Matrices}  
Over time, GPU vendors have added specific hardware features to further accelerate matrix operations, especially for AI. NVIDIA introduced \textbf{Tensor Cores} in its Volta architecture (2017): these are specialized hardware units that perform a fused matrix multiply-and-add on small matrices (e.g., $4 \times 4$ or $8 \times 8$ matrices) in one operation, typically with lower precision like FP16 or INT8. A Tensor Core can, for example, take two $4\times4$ FP16 matrices $A$ and $B$ and an FP16 or FP32 accumulator matrix $C$, and compute $C += A \times B$ in one go. This yields 64 FMAs per clock per Tensor Core. With many Tensor Cores operating in parallel, the throughput is immense – which is why GPUs post such high TOPS (tera-operations) or TFLOPS for matrix arithmetic. These Tensor Core operations correspond exactly to the kind of computation needed in deep learning layers. AMD and other GPU makers have similar matrix-acceleration units in their newer designs as well.

It’s worth noting that before specialized cores, GPUs already had the means to accelerate matrices via normal SIMT parallelism. But Tensor Cores take it further by doing a whole block of multiply-adds in a single instruction, which saves instruction bandwidth and uses dedicated silicon for mixed-precision operations. In essence, it moves the architecture closer to a systolic array (like a mini-TPU inside the GPU, repeated many times). However, from a programmer’s perspective, one usually doesn’t program Tensor Cores directly in assembly; instead, libraries or compilers will use them when possible (for instance, NVIDIA’s cuBLAS or cuDNN will automatically use Tensor Core instructions if the data types and matrix sizes are suitable).

Beyond raw compute, GPUs also organize memory for matrix access. They have coalescing logic to ensure that threads accessing consecutive data in memory (say, thread 0 accesses $A_{0,0}$, thread 1 accesses $A_{0,1}$, etc.) will result in a single combined memory transaction if possible. This favors accessing matrices in row-major order by a warp of threads. For transposed access, either the algorithm rearranges data or uses shared memory to stage it.

\subsection{Example: Matrix Multiply on a GPU}  
To make this concrete, imagine we multiply two $1024 \times 1024$ matrices on an NVIDIA GPU. A common approach (used by algorithms like SGEMM in cuBLAS) is:
- Divide $C$ into tiles of size, say, $128 \times 128$. There will be $8 \times 8 = 64$ such tiles (if we choose 128 for simplicity, though actual tuning chooses tile sizes to fit in shared memory).
- Launch 64 thread blocks, each responsible for computing one $128 \times 128$ tile of $C$.
- Each thread block will further use $128$ threads (organized as, e.g., 16x8 warp arrangement) to compute its tile. It will load a $128 \times k$ slice of $A$ (where $k$ is the inner dimension, here 1024) in portions, and a $k \times 128$ slice of $B$, multiply them, accumulate partial results for its $128\times128$ $C$ tile, and iterate through $k$ in steps (since it can’t load all of $A$ or $B$ at once, it loads, say, $128 \times 16$ sub-tiles incrementally).
- By using fast shared memory, threads in the block collaboratively load sub-blocks of $A$ and $B$, multiply and accumulate before moving to the next sub-block.

This scheme ensures that each element of $A$ and $B$ is loaded a limited number of times (reused within the block) and all ALUs are kept busy. The result is a very high throughput. In fact, “the execution of linear algebra on GPUs is foundational to modern machine learning… The highly parallelized capabilities of GPUs have revolutionized the computational efficiency of many AI operations. Technologies like CUDA and Tensor Cores allow for optimized execution of matrix and vector operations common in ML, significantly accelerating them”【11†L49-L58】. 

In summary, GPUs treat matrices as first-class citizens in terms of performance: their architecture (many parallel cores, high memory bandwidth) and added features (Tensor cores, shared memory, SIMT execution) make them extremely efficient for matrix computations. This is why for matrix-matrix multiplication or large matrix decompositions (LU, QR, etc.), a GPU can outperform a CPU by large factors provided the problem size is big enough to utilize the GPU fully.

\section{Tensors and Tensor Processing Units (TPUs)}  
In mathematics, a \textbf{tensor} is a generalization of scalars (0D), vectors (1D), and matrices (2D) to higher dimensions. For example, a 3D tensor might be an array indexed by $(i,j,k)$. A simple example is a stack of matrices (giving a 3D array), or the kind of multi-dimensional arrays used to represent color images (with dimensions for height, width, and color channels) or mini-batches of data in neural networks (with dimensions for batch size, and features like height, width, channels, etc., making a 4D tensor). Tensors of rank 3 and higher have become increasingly important in modern computing, particularly due to machine learning and data analytics workloads that naturally involve multi-dimensional data.

Operating on higher-rank tensors often involves an operation called \textbf{tensor contraction}, which is a higher-dimensional generalization of matrix multiplication. A tensor contraction is essentially summing (reducing) over one or more indices between tensors. Matrix multiplication, for instance, is a contraction of a 2D tensor with a 2D tensor over one index (the $k$ in the formula $C_{ij} = \sum_k A_{ik} B_{kj}$). More generally, one might have a 3D tensor $X_{i,j,k}$ and a 3D tensor $Y_{k,l,m}$ and contract over the index $k$ to get a 4D result $Z_{i,j,l,m}$, summing over $k$: $Z_{i,j,l,m} = \sum_{k} X_{i,j,k} \times Y_{k,l,m}$. This is just an abstract way of describing a nested loop of multiply-add operations over multi-dimensional arrays.

Hardware like \textbf{Tensor Processing Units (TPUs)} was created to accelerate these kinds of operations, particularly as they appear in deep neural network computations. The Google TPU v1, for instance, was essentially a large matrix multiplier because it turned out neural network inference (for e.g. image classification) could be accomplished mostly with matrix multiplies and accumulations. TPUs are often discussed in terms of tensor ranks because in TensorFlow (Google’s machine learning library), data is represented as n-dimensional tensors and operations are specified as tensor ops. The TPU is optimized for exactly those: matrix multiplies, vector ops, and reductions, which can be used to implement higher-rank tensor operations. In essence, a TPU is a specialized vector/matrix processor where the “vector” is extremely long and implemented via a systolic array feeding a pipeline of multipliers and adders.

Let’s break down the TPU v1 architecture to illustrate how it handles tensors:
- It featured a $256 \times 256$ array of hardware multiply-accumulate units (a systolic array). This can be thought of as 256 processing elements along each dimension, so that it can take in a 256-element chunk of a row of matrix $A$ and a 256-element chunk of a column of matrix $B$ each cycle, multiply corresponding pairs, and sum them along the way. Each cycle, partial results are passed along to the next processing element, like an assembly line (this is the essence of a \textit{systolic array}, where data rhythmically pulses through an array of PEs, each doing a MAC and passing the result on).
- This array is fed by a dedicated on-chip memory (since feeding 65k multipliers with data each cycle requires huge bandwidth). The programming model exposed by Google’s software was that you perform an XLA (Accelerated Linear Algebra) operation which, for example, multiplies two medium-sized matrices at a time using this unit.
- The rest of the TPU core handled things like activation functions, accumulators, and transfers from host memory.

While originally meant for 2D matrix multiplies, the TPU’s systolic array can be utilized for higher-dimension tensor ops by appropriate tiling. For instance, consider a batched matrix multiplication (which is a 3D tensor operation: multiplying two 3D tensors over a shared dimension yields a 3D result). The TPU can process each batch item’s matrix multiply in sequence or partly in parallel if the hardware supports it, but since each matrix multiply is handled by the 2D systolic array, it effectively handles the 3D tensor (batch, i, j) by looping over the batch dimension outside of the systolic array. Similarly, a 4D convolution (common in CNNs) can be lowered to many matrix multiplies (via an \textit{im2col} transformation, which flattens patches of the input feature map and weights into large matrices). TPU and GPU libraries often do exactly that: convert higher-rank tensor operations into a series of matrix operations that the hardware can execute efficiently. This is one of the key strategies in hardware/software co-design: rather than building hardware that directly handles arbitrary high-rank operations, we reduce the problem to the core operations our hardware is extremely good at (matrix multiplies, in the case of TPU/GPU)【21†L61-L69】【21†L63-L71】.

Google’s TPU (and similar “tensor accelerators”) sacrifice generality for performance. As one paper described, “They essentially only multiply matrices, but are very good at doing so”【32†L229-L237】. The TPU can dramatically outperform a CPU or even a GPU on large matrix multiplies, especially at lower precision, because it has a whole 2D array of multipliers running in tandem on the problem. The measured performance of a TPU v3 pod (2048 TPU cores) was about 21 petaFLOPS for FP32 matrix multiplication【32†L231-L239】, which is an astonishing number made possible by this specialized design (for comparison, a contemporary CPU might have ~0.5 TFLOPS of DP throughput per chip, and a high-end GPU on the order of 10-20 TFLOPS of FP32). Of course, that 21 PFLOPS is only achievable on matrix multiplication or very similar operations; a TPU cannot accelerate arbitrary code in the way a CPU/GPU can. This highlights a theme in modern processor design: identifying a key linear algebra operation (like matrix multiply) that dominates certain workloads and building a custom data path to handle it with maximum efficiency.

The term “Tensor Processing Unit” can be a bit misleading, since the hardware deals mainly with matrices. However, the reason for the name is the target application: deep learning models involve tensors (often 4D or higher for things like [batch, width, height, channels] of an image batch, or [batch, sequence\_length, features] for language models, etc.), and the TPU is meant to accelerate the processing of these tensors end-to-end. Since any high-rank tensor operation can be composed (by the compiler) into a series of matrix operations (possibly with reshaping of data in between), the TPU effectively accelerates the tensor operation. In short, the TPU and similar accelerators treat a matrix multiply as the fundamental unit of computation (some have called it the “new transistor” of computing in the ML era). Indeed, a recent architecture called a Tensor Contraction Processor argued that “tensor contraction forms the backbone of machine learning models, handling computations over multi-dimensional data” and that most AI chips focus on matrix multiplication as the basic operation of acceleration【21†L49-L58】【21†L59-L67】. In that paper, they note that typical accelerators map higher-order tensor contractions to sequences of matrix multiplies on fixed-size units, which may not fully exploit data locality or parallelism for irregular tensor shapes【21†L61-L69】. This suggests that future tensor accelerators might incorporate even more flexible ways to handle high-dimensional data beyond just square matrix blocks.

\begin{figure}[hbt]
\centering
\begin{tikzpicture}[>=stealth, node distance=12mm]
\tikzstyle{pe}=[draw, fill=blue!10, minimum size=6mm];
% Define a 2x2 systolic array of PEs
\node[pe] (PE11) at (0,0) {$PE_{11}$};
\node[pe] (PE12) at (8,0) {$PE_{12}$};
\node[pe] (PE21) at (0,-8) {$PE_{21}$};
\node[pe] (PE22) at (8,-8) {$PE_{22}$};
%\node[pe] (PE22) [below=8mm of PE12] {$PE_{22}$};
% Arrows for data flow
\draw[->] ($(PE11)+(-6mm,0)$) -- (PE11.west) node[midway, above] {\small $A$ row};
\draw[->] (PE11.south) -- (PE21.north) node[midway, left] {\small partial sum};
\draw[->] ($(PE12)+(-6mm,0)$) -- (PE12.west) node[midway, above] {\small $A$ row};
\draw[->] (PE12.south) -- (PE22.north) node[midway, right] {\small partial sum};
\draw[->] ($(PE11)+(0,6mm)$) -- (PE12.north) node[midway, above] {\small data to next col};
\draw[->] ($(PE21)+(0,6mm)$) -- (PE22.north) node[midway, above] {\small data to next col};
\draw[->] ($(PE11)+(6mm,0)$) -- ($(PE12)+(-6mm,0)$);
\draw[->] ($(PE21)+(6mm,0)$) -- ($(PE22)+(-6mm,0)$);
\draw[->] ($(PE21)+(-6mm,0)$) -- (PE21.west) node[midway, above] {\small $A$ next row};
\draw[->] ($(PE22)+(-6mm,0)$) -- (PE22.west) node[midway, above] {\small $A$ next row};
\end{tikzpicture}
\caption{Conceptual diagram of a small $2\times2$ systolic array for matrix multiplication. Each Processing Element ($PE_{ij}$) takes an input from the left (elements of matrix $A$) and an input from above (elements of matrix $B$ flowing from a column), multiplies them, and adds to its running sum, then passes the $A$ data to the right and partial sum down. In a larger systolic array like in a TPU, a $256\times256$ grid of such PEs can compute a 256$\times$256 matrix-product in a pipelined wave.}
\label{systolic-fig}
\end{figure}

Figure~\ref{systolic-fig} illustrates a simplified systolic array with 4 PEs, just to give an idea of data flow: the design scales to larger sizes by repeating the pattern. 

The impact of TPUs and tensor accelerators has been significant. They pushed hardware design toward domain-specific accelerators where linear algebra (especially matrix multiply/accumulate) is done orders of magnitude faster than on general CPUs. This is wonderful for deep learning and also has led to exploration of using such units for other domains (scientific computing, graph analytics by linear algebra formulation, etc.). For instance, researchers have tried to run weather simulation or quantum chemistry algorithms on TPUs by expressing them in terms of matrix operations【32†L239-L247】【32†L250-L258】. There are challenges (not every algorithm maps cleanly to matrix multiply, and TPUs were initially designed with only ML in mind, e.g. low precision), but it underscores that much of compute can be seen through a linear algebra lens.

In summary, tensor-oriented processors like TPUs emphasize high-throughput matrix/tensor operations. They handle higher-rank tensors by essentially breaking them into sequences of vector or matrix operations, handled by large parallel units (like systolic arrays). This design achieves extremely high performance for the targeted ops (often at the cost of flexibility). The success in AI has cemented the importance of linear algebra in the future of processor design: rather than just adding wider vectors, companies are adding entire matrix units and tensor accelerators to their chips, which we will see next when examining modern CPU instruction set developments.

\section{Modern Instruction Set Support for Linear Algebra (RISC-V, ARM, x86)}  
We’ve discussed conceptually how different processors handle vectors, matrices, and tensors. Now, we focus on concrete implementations in current processor \textbf{Instruction Set Architectures (ISAs)} and their hardware features. In particular, we look at three prominent architectures: \textbf{RISC-V}, \textbf{ARM}, and \textbf{x86}, in that order, to see how each provides facilities for linear algebra operations. These ISAs each have extensions or features explicitly aimed at accelerating vector/matrix math:

\subsection{RISC-V: Vector Extension and Tensor Proposals}  
RISC-V is a relatively new ISA (introduced in 2014) designed to be modular and extensible. It has a base integer ISA and optional extensions for floating-point, atomic ops, etc. Notably, one of the flagship extensions is the \textbf{RISC-V Vector Extension (RVV)}. The RVV was inspired by classical vector machines (like Cray) and modern needs (like multimedia and ML). It uses a vector-register model where registers (named $v0, v1, \dots$) can hold a flexible number of elements, determined at runtime by a \emph{vector length} setting. Unlike x86/ARM fixed 128-bit or 256-bit registers, RISC-V vectors are \emph{length-agnostic}: a hardware implementation can choose a maximum vector length (say 128 bits, or 512 bits, or 2048 bits), and software can adapt to it. The instructions are written in a way that they operate on $\text{VL}$ elements at a time, where VL is set by an instruction \texttt{vsetvli} based on the data size and hardware max. This approach is sometimes called “LLVM-style vectors” or “vector-length agnostic programming” and is conceptually similar to how ARM’s SVE works as well.

The RISC-V Vector Extension supports a rich set of operations: elementwise arithmetic, dot products, vector reductions (sum, min, etc.), scatter/gather from memory, mask-based predication (operate only on some elements), and even segment operations (for struct-of-array handling). It is a very complete vector ISA, capable of expressing typical linear algebra loops efficiently. For example, RVV has a \texttt{vfmacc} instruction for floating-point fused multiply-accumulate on vectors (accumulating into a vector or scalar), which is great for dot products or matrix multiplies. It also has integer multiply-add with widening (so you can do int8 * int8 accumulate into int32 across a vector, akin to the ARM SDOT). In essence, RVV brings Cray-like vector processing to the modern era with enhancements for software flexibility and integer/FP mix.

One concrete usage: A routine to multiply a matrix by a vector (y = A * x, which involves a dot-product for each row of A) can be vectorized such that each row’s dot product uses the vector unit to do dozens of MACs in parallel. If the vector length (VL) is, say, 256 (meaning 256 operations/cycle), the speedup is enormous versus scalar. The vector unit can also be used to do multiple dot products in parallel (e.g. compute 8 partial sums simultaneously).

Beyond the vector extension, the RISC-V community and companies working with RISC-V have proposed even more specialized \textbf{matrix extensions}. For instance, Alibaba’s T-Head (Xuantie) has defined a Matrix Multiply Engine (MME) extension that adds new instructions and registers specifically for matrix tiles, aimed at accelerating small matrix multiplications common in AI. This extension allows loading entire small matrices with one instruction and performing matrix-multiply with another, similar to how Intel’s AMX (see below) works. A recent academic proposal “MX” (Matrix extension) showed that adding just a small tile buffer to an RVV implementation (to hold sub-blocks for reuse) can significantly boost energy-efficiency of matrix multiplication【12†L57-L65】【12†L67-L70】. There’s also the concept of a \textbf{Tensor Unit} as a peripheral: for example, Semidynamics (a RISC-V core company) announced a “world’s first fully-coherent RISC-V Tensor Unit” which attaches to their vector unit【16†L139-L147】【16†L149-L157】. This Tensor Unit essentially acts as a matrix-multiply accelerator (systolic-like) tied into the RISC-V core’s cache coherence, so software can offload matrix ops to it seamlessly. They report a 128x speedup on AI workloads versus just running on the scalar core【16†L149-L157】. The integration with the vector register file means the tensor unit can pull matrices from the vector registers, do a fast multiply, and deposit results back, which is a clean design.

In summary, RISC-V’s strategy is to provide a very powerful and flexible vector ISA for general use (which covers HPC and many linear algebra needs), and allow specific implementations to add even more acceleration (matrix/tensor units) as optional extensions. The base vector ISA (RVV) is quite capable: it can, for example, implement a 64x64 matrix multiplication in software using vector instructions to multiply-add rows and columns with fairly high efficiency【13†L73-L80】. But for those needing even more, the ecosystem is exploring direct matrix operations. Since RISC-V is open, these extensions are being discussed openly, and perhaps in the future an official matrix extension will be standardized.

\subsection{ARM: NEON, SVE, and SME (Scalable Matrix Extension)}  
ARM architecture has evolved to include increasingly robust linear algebra support. ARM’s first significant step was \textbf{NEON}, introduced around ARMv7 (for 32-bit) and integral to ARMv8 (64-bit) cores. NEON is a SIMD vector extension with 128-bit vector registers, capable of handling typical data types (8,16,32-bit integers, 32- and 64-bit floats). It was primarily intended for media (audio/video) and signal processing, so it includes multiply-accumulate, wide loads, etc. Many mobile workloads (like image processing, simple ML filters) use NEON to speed up linear algebra kernels (e.g., small matrix multiplies for audio filters or image convolutions). NEON can do 4 FP32 operations per instruction (since 128 bits/32 = 4), or 8 FP16, or 16 8-bit operations, etc. This is analogous to Intel SSE.

For larger-scale vector needs, ARM introduced \textbf{SVE (Scalable Vector Extension)} with ARMv8.2-A (implemented notably in Fujitsu’s A64FX CPU that powers the Fugaku supercomputer). SVE is conceptually similar to RISC-V’s vector extension: the actual vector register width is not fixed in the ISA (it could be 128b, 256b, … up to 2048b), and code is written in a vector-length-agnostic way. SVE adds powerful predication (per-element masks) and gathers/scatters, making it very suitable for HPC workloads that involve vectors and matrices with irregular memory access. For linear algebra, SVE can apply operations on very long vectors, which benefits routines like DGEMM (dense matrix multiply) by enabling operations on many elements at once and by helping compilers auto-vectorize loops without knowing the exact hardware width.

One interesting addition in ARMv8.4-A was the introduction of \textbf{dot product instructions} in NEON【44†L17-L25】. These instructions (like SDOT, UDOT) take four 8-bit elements from one vector and four 8-bit from another, multiply pairwise and accumulate into a 32-bit result (and actually they do this for four such groups in one 128-bit vector operation, effectively computing 16 multiply-accumulates per instruction). This was clearly motivated by neural network inference, allowing 8-bit multiplication with accumulation into 32-bit (to avoid overflow) – a common pattern in conv and fully-connected layers. By using SDOT/UDOT, an ARM CPU core could significantly speed up int8 matrix multiplication or convolution compared to doing it element by element. For example, multiplying two $4\times4$ int8 matrices and accumulating could be done with a handful of these dot instructions. In fact, ARM’s documentation showed over 17% performance improvement in certain video encode workloads at 1080p just by utilizing UDOT instructions in NEON【44†L1-L9】【44†L11-L19】. This is a case where a linear algebra primitive (small dot products) was added to a general ISA to accelerate a specific domain (ML and DSP).

Moving to ARMv9, ARM introduced the \textbf{SME (Scalable Matrix Extension)}. SME builds on SVE to provide even more acceleration for matrix operations. In essence, SME adds a new set of registers (the “ZA” matrix storage) that can hold two-dimensional blocks of data, and instructions that operate on those blocks performing matrix multiplies efficiently. It’s ARM’s counterpart to things like Intel’s AMX (discussed below). SME is scalable, meaning it is designed to work with different matrix sizes and tile shapes, and it is particularly aimed at AI workloads. An ARM overview states: “SME is an ISA extension in Armv9-A which accelerates AI and ML workloads and enables improved performance for matrix operations such as matrix multiplication”【29†L5-L13】【29†L23-L31】. Essentially, SME allows the CPU to handle operations on, say, 16x16 or 32x8 blocks of matrices with single instructions, offloading loops from the software into hardware. The details of SME include:
- A new matrix register file ZA (which can be viewed as a collection of vector registers grouped to form matrix tiles).
- New instructions to initialize, load, and compute on these ZA matrices. For example, one can load a tile of a matrix with one instruction (instead of many strided vector loads), and then perform a matrix multiply accumulate of a tile of A with a tile of B, accumulating into a tile of C.
- It supports transposition and other addressing modes convenient for matrix math.

SME also retains the “scalable” nature: an implementation can choose how big the physical matrices are that it can operate on at once. It might be 128-bit wide vectors combined into a certain tile shape. The key benefit is reducing the overhead of issuing lots of individual MACs and instead doing a whole block in fewer instructions.

To summarize ARM’s approach: basic SIMD with NEON for small vectors, scalable SIMD with SVE for large vectors and easier auto-vectorization, and now SME to tackle matrices directly. This hierarchy means ARM cores are becoming more and more like linear algebra engines:
- A programmer (or compiler) can use NEON intrinsics to speed up say a $4\times4$ matrix multiply with dot instructions.
- Or use SVE to vectorize a loop that multiplies a $1000\times1000$ matrix by a vector.
- Or use SME to perform a $64\times64$ matrix multiply in a few steps by loading tiles and computing.

ARM’s development was clearly driven by both HPC (e.g., Fujitsu’s interest in vectorizing for supercomputing) and AI (everyone’s interest in accelerating ML). The inclusion of these in a general-purpose CPU means that, e.g., a smartphone or laptop with an ARMv9 chip will have significant linear algebra capabilities even without a GPU, which can be useful for moderately sized ML tasks or heavy signal processing.

\subsection{x86: SIMD to AVX-512 and AMX}  
The x86 architecture (used by Intel and AMD) has seen a long evolution of SIMD and matrix support:
- In 1997 Intel introduced \textbf{MMX} (64-bit registers, integer ops mainly for multimedia).
- In 1999, \textbf{SSE} added 128-bit registers and support for single-precision floats, then double-precision in SSE2, and so on. SSE enabled doing four 32-bit float ops at once, which was a boon for graphics and basic linear algebra on PCs.
- Around 2011, \textbf{AVX} (Advanced Vector Extensions) expanded SIMD registers to 256 bits (YMM registers). AVX and AVX2 allowed doing eight 32-bit floats or four 64-bit doubles in one instruction, etc. This was a significant boost for linear algebra code on CPUs; libraries like Intel’s MKL and OpenBLAS could now utilize 8-wide vector units, doubling the throughput compared to SSE.
- \textbf{FMA3/FMA4} extensions also came (FMA3 in Intel Haswell, 2013) which provided the fused multiply-add instruction on these wide registers (e.g., do 8 parallel FMAs on 8 floats in one YMM register with one instruction). This directly improves matrix multiply kernels, as the inner loops can use FMA to do multiply-accumulate in one step.
- In 2015, \textbf{AVX-512} was introduced (Knights Landing, then Skylake-SP for general servers) with 512-bit ZMM registers. AVX-512 can do 16 single-precision floats or 8 doubles in one vector operation, and added comprehensive masking (like predication for vectors) and many new instructions (integer vector scatter/gather, etc.). This was another doubling of vector width and a leap in capability; a single AVX-512 FMA can perform 16 float multiplications and additions in one go, which is quite powerful. The Linpack (HPL) benchmark which measures dense linear algebra performance on CPUs saw big jumps with AVX-512 due to increased parallelism【30†L17-L26】【30†L23-L28】.

However, x86 did not stop at wide SIMD. Inspired by the needs of deep learning, Intel introduced a feature called \textbf{AMX (Advanced Matrix Extensions)} in its 4th-generation Xeon Scalable (codename Sapphire Rapids, released 2022). AMX is a new paradigm: instead of just widening vectors further, it provides a 2D register file (called “tiles”) and operations on these tiles, focusing on matrix multiply-accumulate for int8 and bfloat16. The AMX provides a set of tile registers (eight registers each capable of holding, for example, 16x64 int8 values or 16x32 bfloat16 values), and a single instruction can multiply two tiles (one treated as 16x64 and another as 64x16, producing a 16x16 result) and accumulate into a third tile. In total, one such operation computes 16*16*64 = 16384 integer multiply-accumulates in one instruction【29†L13-L18】 (if using int8). This is an enormous step up in density of computation per instruction. Essentially, it’s like having a fixed-size matrix multiplication accelerator inside the CPU. The idea is to accelerate inner kernels of deep learning (e.g., matrix multiplies in transformers and CNNs) to be competitive with specialized AI accelerators. One can imagine the programming model: you load a block of matrix A into a tile, a block of B into another, execute the tile matrix multiply instruction, and then move the result out or continue accumulating. This offloads what would be a triple nested loop into microcode/hardware.

It’s worth noting that typical x86 CPUs also support lots of accumulative operations that help with linear algebra:
- They have reduction operations like \texttt{DPPS} (dot product of packed singles) which was introduced with SSE4. This computes a dot product of two vectors of 4 floats, giving a single float result (with some masking control). This is a limited case, but was useful in graphics shaders and linear algebra kernels to avoid having to do a horizontal add manually.
- There are also instructions for permutations, transposition of 4x4 matrix in registers (useful for data layout adjustments), etc., which are used by libraries to optimize memory access patterns for matrix ops.

AMD’s x86 implementations have similar support (though AMD did not implement AVX-512 in its Zen cores as of 2023, focusing on 256-bit AVX2, they may adopt something similar to AMX in future or rely on GPU for big matrix tasks).

In summary, x86 started with basic SIMD for vectors but has grown to include very powerful vector units and now even matrix-processing capabilities. A programmer using x86 could utilize intrinsics or libraries to:
- Use AVX-512 to do linear algebra in parallel (for double precision heavy HPC, AVX-512 is typically the go-to).
- Use the VNNI (Vector Neural Network Instructions, an AVX-512 subset) which introduced an instruction to do multiply-add of int8 with automatic accumulation in int32, effectively like a dot-product across 4 elements (similar to ARM’s dot, but across more elements because of wider registers)【11†L53-L57】.
- Use AMX for very high throughput int8 or bfloat16 matrix ops, which would be especially useful for inference workloads. 

All these are accessible to developers through libraries (like oneDNN, Intel’s optimized DNN library, will use AMX if available) or assembly if doing custom kernels. It represents the convergence of CPU design with the needs of linear algebra: where once one might offload to a GPU or TPU for massive matrix multiplies, we now see CPUs themselves incorporating mini-TPUs internally (for specific data types) to handle those tasks.

One should note that while these advanced features exist, using them efficiently requires careful programming or compiler support. Automatic vectorization by compilers can target SSE/AVX for typical loops, but making use of AMX or even fully utilizing AVX-512 often requires hand-tuned code or high-level libraries. Nonetheless, the hardware capability is there and is a testament to how important linear algebra has become even in general-purpose processors.

Table~\ref{isa-table} summarizes some of the key linear algebra-related features of RISC-V, ARM, and x86:

\begin{table}[hbt]
\centering
\begin{tabular}{|l|c|c|c|}
\hline 
\textbf{Feature} & \textbf{RISC-V} & \textbf{ARM} & \textbf{x86 (Intel/AMD)} \\
\hline 
Baseline SIMD & None in base; & NEON (128-bit) & SSE (128-bit) \\
 & Vector Ext (VLEN config.) &  & AVX (256-bit), AVX-512 (512-bit) \\
\hline
Max vector width & Flexible (impl-defined, & 128-bit NEON; & 512-bit (AVX-512) \\
 (standard) & e.g. 128, 256, 512 bits) &  up to 2048-bit SVE & (fixed width) \\
\hline
Notable ops & Vector FMA, reductions, & NEON: FMA, SDOT/UDOT (8-bit dot); & AVX-512: FMA, VNNI (8-bit MAC); \\
 & mask/predication & SVE: predication, gathers; & horizontal add, permutes; \\
 &  & SME: tile load, matmul & AMX: tile matmul int8/bf16 \\
\hline
Special matrix/tensor & Proposed Matrix Ext, & SME (Matrix Extension in v9); & AMX (Tile matrix mult. unit); \\
 support & Tensor Unit (impl.)【16†L139-L147】 & also GPUs in SoC for AI & GPUs in discrete/SoC for AI \\
\hline
\end{tabular}
\caption{Linear Algebra support in modern ISAs. RISC-V relies on an advanced Vector extension and allows custom matrix units. ARM provides NEON for basic SIMD, SVE for scalable vectors, and SME for dedicated matrix ops. x86 started with SSE/AVX vectors and now includes AMX for matrix multiplication acceleration.}
\label{isa-table}
\end{table}

The above table is a high-level comparison. Each platform has nuances (e.g., rounding modes for FMA, whether integer ops saturate or wrap, etc.) that we skip here. But clearly, all major ISAs have recognized the importance of vectors and matrices, adding instructions or coprocessors to accelerate them. In the end, whether it’s called a “vector extension” or a “matrix extension”, the hardware aim is the same: do more multiply-adds in parallel, because that’s what so many linear algebra problems demand.

\section{Mapping Linear Algebra Operations to Hardware}  
We have talked about the hardware and ISA features individually; now let’s consider how typical linear algebra operations map onto these hardware features. The goal is to understand, for a given mathematical operation, what the processor is doing under the hood and which part of the hardware is leveraged.

\subsection{Dot Product (Vector-Vector) Mapping}  
A dot product $s = x\cdot y = \sum_i x_i y_i$ is conceptually a single number but involves $n$ multiplications and additions. On a modern processor:
- A simple implementation would be a loop that multiplies one pair at a time and accumulates into $s$. This would use the scalar ALU or FMA unit $n$ times.
- With SIMD, the processor can do, say, 8 multiplications in parallel with an instruction and then another to add them to an accumulator vector. But you also need to add across the vector lanes to get a single sum (a horizontal reduction). For example, using AVX2, one might do 8 multiplies (producing 8 partial products), then use an instruction like \texttt{VPADDD} or \texttt{VHADDPS} to halve them to 4, then 2, then 1. Some ISAs provide a direct “vector reduce” instruction: RISC-V has \texttt{vredsum} as used in our earlier pseudocode, which sums all elements of a vector into a scalar【13†L37-L45】【13†L73-L80】. Lacking that, compilers emit a tree of adds or rely on special dot instructions.
- On architectures with a dot-product instruction (like ARM’s SDOT), the dot of small segments can be done directly. E.g., SDOT computes 4 products + sum in one go (for int8). So the loop can accumulate 4 at a time into each 32-bit lane.
- On a GPU, each thread might handle one dot product (if the vectors are large, you could also split one dot product across multiple threads, but usually a thread or warp per dot is fine). Each thread then uses its GPU core’s FMA units in a loop. However, GPUs often use each thread to compute one element of a larger result (like a matrix multiplication), rather than just a single dot product in isolation.
- Specialized hardware: DSPs or older vector machines might have had a specific accumulator register that is wider to hold a sum without overflow, and a MAC unit feeding it every cycle. For instance, some DSPs have 40-bit accumulators for 16-bit data dot products to avoid overflow.

Thus, a dot product maps to the general feature of “fused multiply-add repeated with accumulation”. All modern hardware has something to support that, whether it’s a scalar FMA unit used in a loop or a full SIMD FMA used in a vectorized loop. The main challenge is summing up results without losing precision (for floating point, order of summation matters). Some architectures (like x86 with x87 in old days) had 80-bit extended precision to reduce error in sums; nowadays algorithms sometimes do pairwise summation or use Kahan summation in software if needed for very high accuracy, since the hardware will do a straightforward sum.

\subsection{Matrix Multiplication Mapping}  
Matrix multiplication $C = A \times B$ (for simplicity, assume all square $n \times n$ here) is typically implemented in a blocked manner on CPUs. The reason is memory hierarchy: a naïve triple loop will make $n^3$ accesses and $n^3$ operations; we want to reuse data in caches to reduce memory traffic. The classical approach is known as \textbf{blocked matrix multiplication}: break matrices into submatrices (blocks) that fit in cache, then multiply those blocks using the same algorithm recursively or with loops.

On hardware with SIMD:
- The innermost part that does a rank-k update of a submatrix (like $C_{sub} += A_{sub} \times B_{sub}$ where these are, say, $r \times k$ times $k \times c$ matrices) can be vectorized. Typically one would vectorize the multiplication of a row of $A_{sub}$ by a column of $B_{sub}$ (which is a dot product). Many libraries implement matrix multiply by computing several elements of C in one inner loop to reuse A and B. For example, an optimized kernel might compute a $4\times4$ block of C in registers: it will load 4 elements of a row of A into a vector register (with replication to pair with 4 different columns of B), and 4 elements of a column of B into another vector, perform 4 FMAs (to update 4 accumulators for 4 different C elements) – effectively computing 4 dot products in parallel. This is sometimes called the “outer product” method: take one row of A (as a vector) and one column of B (as a vector), outer product to update a block of C. Contrast with the “inner product” method: computing one C element at a time as a dot of a row and a column.
- Modern CPU codes often use intrinsics or assembly to unroll and vectorize such operations. For instance, using AVX, one might load 8 elements of A (two 4-length vectors if doing double precision, or one 8-length if single precision) and broadcast each to multiply with a vector from B. AVX-512 introduced new instructions like \texttt{VBROADCASTSS/VBROADCASTSD} to efficiently broadcast one element to all lanes for such usage. It also introduced \texttt{VPDPBUSD} (Packed Dot Product for Byte and Unsigned Byte with Saturation?) in AVX-512-VNNI, which does dot products of 8-bit integers with accumulation – again, targeted at ML int8.
- On GPUs, as we described, the mapping is typically one thread block per submatrix. Each thread in the block might compute one element of C or one small tile of C. They use shared memory to hold tiles of A and B. This is effectively a software-managed cache, because early GPUs didn’t have big caches; now they do have caches but shared memory (scratchpad) is still explicitly used for control. The GPU cores execute the inner multiply-add loops very efficiently since each thread can issue an FMA every cycle and with so many threads, you achieve a high degree of ILP (instruction-level parallelism) and TLP (thread-level parallelism).
- TPUs map this to systolic array hardware. The mapping is more direct: if the matrix is larger than $256\times256$, it’s done in blocks that size. If smaller, not all PEs are used. The compiler (XLA) schedules the multiply-add flow so that the data is loaded and fed into the systolic array at the right times.
- With new matrix ISA support like ARM’s SME or x86’s AMX, the mapping becomes even more straightforward: load tiles and call the tile matrix multiply instruction. For example, on an AMX-equipped CPU, an implementation of matrix multiply could:
  1. Set up the tile configuration (tile size).
  2. Loop over $k$ dimension in steps of tile width (e.g., 64 if tiles are 16x64).
  3. For each step, load a 16x64 tile from A (16 rows, 64 cols) into tile register 0, load a 64x16 tile from B (64 rows, 16 cols) into tile register 1.
  4. Execute \texttt{TPBFMULD} (just a hypothetical name) to multiply these into a 16x16 tile accumulator for C.
  5. After looping, the tile accumulator holds the 16x16 block of C, which can then be stored to memory.
  
  In this way, a lot of the index math and innermost multiplies are offloaded to the hardware performing the tile operation. The effect is similar to what well-optimized assembly would do, but now the CPU has explicit support to make it easier and faster.

As mentioned, \textbf{tensor contraction} is a generalization of matrix multiply. In many cases, a general tensor contraction (with more than two input dimensions) is implemented by rearranging the data (if necessary) and calling a matrix multiply. Compilers like LLVM, or domain-specific compilers for ML, often use a tactic where they identify a tensor contraction (say using Einstein summation notation) and then choose one index to act as the “summation index” and treat the rest as combined outer indices, converting it to a matrix multiplication【21†L53-L61】【21†L63-L67】. For example, consider a contraction $Z_{a,b,d,e} = \sum_{c} X_{a,b,c,d} * Y_{c,e}$. Here $X$ is a 4D tensor and $Y$ is 2D, contracting over index $c$. We can reinterpret $X$ as an $(a*b*d) \times c$ matrix (by flattening the first and fourth dims with $a, b, d$ combined as one dimension of size $a*b*d$ and $c$ as the other dimension), and $Y$ as a $c \times e$ matrix, and the result $Z$ as an $(a*b*d) \times e$ matrix which can then be reshaped to $(a,b,d,e)$. This kind of reshaping and using GEMM is extremely common in high-performance libraries because highly tuned GEMM (matrix multiply) routines exist, so any tensor operation is ideally funneled into GEMM calls if possible, a strategy sometimes called “GEMMification”. It leverages the fact that hardware (and low-level software) for matrix-multiply is very optimized【21†L63-L71】【21†L73-L78】.

When direct mapping to hardware isn’t possible (for example, a very irregular contraction), compilers or libraries will break the loops and use the available primitives as much as they can, perhaps handling some part in software. But given the trend, more and more types of contractions are being handled in hardware-friendly ways.

One important consideration is \textbf{data movement and memory layout}. Hardware features like caches, TLBs, etc., strongly influence performance of large linear algebra operations. That’s why techniques like blocking, tiling, using transpose to align memory accesses, etc., are crucial. Software often will pad matrices or choose strides carefully to get good alignment for SIMD and avoid cache conflicts. Modern compilers use the \textit{loop nest optimization} (tiling, unrolling, etc.) to maximize the use of the CPU’s vector units and caches for matrix operations【30†L39-L47】【30†L51-L58】. For example, the classic approach to matrix multiply on CPU might be 3-level nested loops $i,j,k$, and a smart compiler or library will tile it into four nested loops: $i$ in steps of $I$, $j$ in steps of $J$, $k$ in steps of $K$, and inside do an $I\times J \times K$ micro-kernel that fits in L1 cache and uses SIMD registers for $K$ loop unrolling. Libraries like BLIS or MKL have this down to a science, often achieving near-peak FLOPS of the CPU (e.g., $>90\%$ of the theoretical peak FLOPS for large matrix multiply)【30†L25-L33】.

\subsection{Higher-Rank Tensor Decomposition and Mapping}  
When dealing with tensors of rank $>2$, one strategy beyond just turning them into matrices is to use mathematical \textbf{tensor decomposition} techniques. These are not exactly hardware features, but rather algorithmic transformations that can drastically reduce computation if the tensor has certain low-rank structure. There are several popular tensor decompositions: 
- **CP (CANDECOMP/PARAFAC) Decomposition**: expresses a tensor as a sum of outer products of vectors. For a 3D tensor $T_{i,j,k}$, a rank-$R$ CP decomposition writes $T_{i,j,k} \approx \sum_{r=1}^R a_{i,r}\,b_{j,r}\,c_{k,r}$. If $R$ is much smaller than the typical dimension, this is a big compression. Computing with the decomposed form can be faster: e.g. multiplying this tensor by something might be easier since you deal with the smaller factors $a, b, c$.
- **Tucker Decomposition**: expresses a tensor as $T_{i,j,k} \approx \sum_{p,q,r} G_{p,q,r}\, A_{i,p}\, B_{j,q}\, C_{k,r}$ with a smaller core tensor $G$. This is like a higher-order SVD.
- **Tensor Train (TT) or Matrix Product State**: factorizes a high-D tensor into a chain of 3D tensors which are multiplied together.
- **Tensor Ring**: similar to TT but in a loop.

These decompositions are used to compress neural network models (to reduce parameter count and compute)【18†L45-L54】【18†L51-L59】. For instance, a fully-connected layer weight matrix (which is 2D) might be reshaped to a 4D tensor and then decomposed by Tucker or TT to yield a sequence of smaller matrices to multiply, reducing ops. Or a convolutional kernel tensor (4D: out\_channels, in\_channels, filter\_height, filter\_width) can be decomposed into, say, a sequence of smaller 4D tensors or a CP decomposition (sum of outer products of four vectors). There have been successful experiments showing significant compression with minor accuracy loss, and these reduce the number of MAC operations needed【17†L5-L13】【17†L23-L30】.

From a hardware perspective, leveraging these means:
- The software (framework or library) detects a possibility to use a decomposition, or the model is trained with decomposition (some research trains networks with low-rank constraints).
- Then instead of one big matrix multiply, the program will perform multiple smaller ones. Those smaller ones fit better in caches or can make use of limited precision better. Sometimes it also allows parallelizing in a different way.
- Example: Suppose a weight matrix $W$ of size $1000 \times 1000$ is low rank $R=50$. We can factor $W = P \times Q$ where $P$ is $1000\times 50$ and $Q$ is $50 \times 1000$. Then $Wx = P \times (Qx)$: first compute $u = Qx$ (which is 50x1000 by 1000x1, cheaper than 1000x1000), then $y = P u$ (1000x50 by 50x1). This does $1000*50 + 1000*50 = 100,000$ multiplies instead of $1000*1000 = 1,000,000$ – a 10x reduction in ops (plus some overhead for addition if rank was approximate). On hardware, this means we perform two smaller GEMMs, which could also be more cache-friendly (less data to move at once).
- Another example: In CNNs, a 3x3 conv on 64 channels can be seen as a 9x64x64 tensor (9 positions in the kernel, 64 input maps, 64 output maps). CP decomposition could approximate this 3D tensor as a sum of outer products of three vectors: one of length 9, one of length 64 (for input), one of length 64 (for output). That effectively breaks the conv into 3 separate convs (or 3 FC layers) with smaller weight vectors. Each of those is far cheaper than the original convolution. On hardware, you might implement these with three smaller matrix multiplications (or three smaller conv calls).

The challenge is that not every tensor is low-rank; sometimes the rank needed for good accuracy is not much smaller than the full size, so no gain. But many real-world tensors do have redundancy.

From the hardware view, we are basically trading one large linear algebra op for a sequence of smaller ones. If the total arithmetic is less, we win on compute. But note, performing multiple passes might increase memory traffic (because we might have to read the input multiple times for each factor). The ideal scenario is when the decomposition dramatically cuts compute and the hardware was compute-bound (like on TPU or GPU which often are compute-bound in dense ops). If memory-bound, a decomposition might not help unless it also cuts memory (which model compression does, by storing fewer weights).

Modern ML frameworks can incorporate these ideas. Some research compilers even automatically factor tensors or use mixed sparse/dense representations (structured sparsity can be seen as a form of decomposition too). Hardware support for this is indirect: by having fast small matrix multiplies, the overhead of doing multiple multiplies is low. For example, if a decomposition says do 10 small matmuls instead of 1 big one, a CPU with good cache might handle that well if those small ones fit in L1/LLC, whereas one big one might spill to DRAM a lot. There’s active research on hardware-aware tensor decomposition【17†L1-L9】【17†L23-L31】, trying to find decompositions that map well to actual hardware (taking into account vector units, memory, etc.)【18†L51-L59】【19†L1-L4】.

In summary, higher-rank tensors can often be flattened or factorized to lower-rank forms, and doing so usually means turning the problem into a set of vector or matrix operations that current hardware is very good at. This is another aspect of co-design: sometimes we modify the problem (via math) to fit the hardware rather than designing hardware to fit the abstract problem. We’ll see more of this in software frameworks.

\section{Software and Hardware Co-Design for Linear Algebra}  
The synergy between software and hardware is critical in extracting performance from linear algebra operations. Software – including low-level libraries, compilers, and high-level frameworks – has evolved in tandem with hardware capabilities. In this section, we explore how software utilizes the hardware features discussed, and how hardware design has been influenced by software demands.

\subsection{Linear Algebra Libraries and BLAS}  
One of the earliest examples of co-design is the development of the \textbf{BLAS (Basic Linear Algebra Subprograms)} standard and libraries. BLAS defines a set of routines for vector and matrix operations (Level 1 BLAS: vector-vector, Level 2: matrix-vector, Level 3: matrix-matrix). By coding algorithms (like solving $Ax=b$ or computing eigenvalues) in terms of BLAS, one can rely on highly optimized implementations of BLAS for each architecture. Library developers (like those behind IBM ESSL, Intel MKL, OpenBLAS, ATLAS, etc.) hand-tune BLAS assembly to make the best use of the machine: using appropriate instruction scheduling, prefetching, SIMD instructions, etc. For instance, the BLAS routine SGEMM (single precision matrix multiply) has been endlessly optimized on each new CPU generation to use wider SIMD or new FMAs. The result is that high-level software (like NumPy, MATLAB, or SciPy) just calls BLAS and automatically gets near-peak performance without knowing the hardware details.

It’s a symbiotic relationship: hardware designers know that if they add a new capability (say AVX-512), library writers will update BLAS to exploit it, improving performance of all applications that rely on BLAS. Conversely, the popularity of BLAS has sometimes influenced hardware; e.g., the presence of a fused multiply-add in hardware is extremely beneficial for BLAS Level 3 routines (matrix multiply)【30†L17-L25】【30†L25-L33】. In HPC, benchmarks like LINPACK (which is basically a measure of how fast a system solves dense linear systems using LU factorization, which is built on GEMM) drive vendors to optimize these operations heavily. The so-called “DGEMM inner-kernel” was the focus of a lot of innovation in the 2000s, as it determined LINPACK scores. This led to improvements in cache hierarchies and prefetching mechanisms in CPUs (to better feed the multiply pipelines).

\subsection{Compiler Optimizations}  
General-purpose compilers (GCC, Clang, ICC, etc.) have evolved to automatically optimize loops that look like linear algebra. Auto-vectorization is a feature where the compiler detects that a loop can be executed with SIMD instructions and emits them. For example, a loop `for(i=0;i<n;i++) C[i]=A[i]+B[i];` will typically be auto-vectorized to use AVX or NEON instructions to add multiple elements per iteration. Compilers also unroll and tile nested loops for matrix operations if they detect a pattern. However, auto-vectorizers have limitations and often don’t rival hand-written assembly for complex operations like matrix multiply, which is why BLAS libraries still reign for those.

There are also domain-specific compilers and languages for linear algebra. MATLAB or Julia, for instance, have JIT compilers that try to fuse operations to reduce temporary memory usage. For example, if you write `z = A*x + y` (with $A$ matrix, $x,y$ vectors), a smart environment will recognize this as a single matrix-vector multiply plus vector add, and if possible, call a single fused kernel rather than doing `A*x` then adding y separately (which would create an intermediate result). This reduces memory traffic and can use fused instructions.

Another aspect of co-design is using \textbf{compiler intrinsics} or built-ins to manually guide vectorization. Most ISAs provide intrinsics (C/C++ functions that map directly to specific instructions, like \texttt{\_mm256\_loadu\_ps} for AVX load, \texttt{\_mm256\_fmadd\_ps} for FMA). Library developers use these to write portable yet optimized code targeting SIMD without writing raw assembly. The compiler then handles register allocation etc. but the developer has fine-grained control of instruction usage.

For high-level tensor operations, compilers like \textbf{XLA (Accelerated Linear Algebra)} from Google take a whole computation graph (like a neural network graph) and perform global optimizations, like combining operations, choosing memory layouts, etc., before lowering to hardware instructions【31†L15-L23】【31†L25-L33】. XLA, for example, might transform a series of tensor ops into one large fused kernel to run on TPU or GPU, to minimize overhead. It treats linear algebra operations as first-class optimizable things. In fact, XLA is part of a broader trend of \textbf{MLIR (Multi-Level IR)} and other such projects to have compiler infrastructure that understands linear algebra at an abstract level (operations like tensor contraction, convolution, etc.) and can then generate low-level code for different backends (CPU, GPU, TPU) efficiently.

It’s worth noting that \textbf{Graphical algorithm mapping} is also used: e.g., convolution algorithms (like Winograd’s algorithm) reduce the number of multiplications at the cost of a few more additions, and compilers/hardware co-design might decide to use such algorithms if they are beneficial on a given hardware. For instance, some GPU libraries use Winograd for certain small conv filters as it reduces math, but if the hardware has excess FMA capacity and is memory-bound, they might stick to direct convolution.

\subsection{Frameworks and Run-Times}  
In the realm of machine learning, frameworks like TensorFlow, PyTorch, MXNet, etc., play a big co-design role. These frameworks typically:
- Express computations in a high-level dataflow (e.g., matrix multiplies, activations, etc. as nodes).
- Use optimized kernels (often from BLAS or specialized libraries like cuDNN for conv, etc.) for each node where possible.
- Or use JIT compilers (like PyTorch’s Glow, or TensorFlow XLA) to generate efficient fused code.

The frameworks also handle selecting the right hardware. For instance, if running on CPU, they may use oneDNN (Intel’s oneAPI Deep Neural Network library, formerly MKL-DNN) which uses AVX/AMX. On GPU, they use cuBLAS or cutlass kernels from NVIDIA. On TPU, they invoke XLA compiled code.

The term \textbf{software/hardware co-design} is very explicit in projects like Google’s TPU: the XLA compiler was built alongside the TPU hardware. Certain design decisions in TPU (like having a large unified accumulator or using bfloat16 format) were made to simplify the compiler and usage. The software knows the TPU does matrix ops of a certain max size, so it will split any larger ones accordingly. Conversely, the TPU’s programming model was designed to fit typical neural network layer patterns so that the compiler could easily place operations on the systolic array without complicated scheduling from the programmer’s side【31†L19-L27】【32†L227-L234】.

Another example is \textbf{Halide} (a domain-specific language for image processing, which is essentially linear algebra on images). Halide allows the programmer to define an algorithm (like a convolution or an optical flow computation) and then separately define a \textit{schedule} which specifies how to tile, vectorize, unroll, etc. Halide’s compiler then emits optimized C++/vector code for CPU or CUDA for GPU accordingly. This separation of algorithm from optimization enables rapid co-design: one can change the schedule (which is like an abstraction of hardware mapping) to suit a new CPU or GPU without changing the algorithm logic.

At a higher level, one could consider languages like Python (with NumPy) or MATLAB as part of co-design: they provide an easy way for users to specify linear algebra operations, which under the hood call optimized library routines. The user doesn’t manually loop and vectorize (which is error-prone and not portable); they just call, say, `C = A.dot(B)` and the environment ensures that uses the best method on the available hardware (be it multi-threaded CPU BLAS, GPU cuBLAS, or even distributed computing library if $A$ is distributed).

\subsection{Parallel and Distributed Linear Algebra}  
Modern supercomputers involve distributed memory (clusters of nodes). Libraries like ScaLAPACK or distributed TensorFlow are designed to split large matrices across nodes and use network communication along with local BLAS to perform linear algebra at scale. Co-design at this level involves the network interface hardware too. Technologies like NVIDIA’s NVLink and Infiniband with RDMA are tuned to quickly shuffle matrix blocks between GPUs or between nodes so that a large matrix multiply can be split, computed in parallel, and results combined. For instance, the SUMMA algorithm for distributed matrix multiply is often used, and hardware features like offloading collective operations or using smart NICs can accelerate the necessary communication. Google’s TPU pods likewise have high-speed interconnect to share results of sub-matrix multiplies between chips【32†L231-L239】【32†L239-L247】.

So co-design extends to considering the entire system: CPU/GPU/TPU with their vector/matrix units, plus the memory subsystem, plus the networking between processors, plus the software algorithms that orchestrate the work.

To highlight co-design implications in a current context: consider \textbf{Transformer models} (used in language AI). They involve very large matrix multiplies (for attention and feed-forward layers). NVIDIA, Google, etc., ensure their libraries and hardware support mixed precision (FP16 or bfloat16) with accumulation in higher precision, which gives performance and sufficient accuracy. Software automatically casts to these types and calls the tensor core kernels. Moreover, frameworks use techniques like \textbf{quantization} (INT8 inference) and the CPU/GPU provide corresponding instructions (like ARM’s dot, x86’s VNNI) to make quantized linear algebra fast. This is a co-design win: you reduce precision to reduce compute, but then you rely on hardware support to not lose too much performance when handling 8-bit values (since doing 8-bit arithmetic on a 32-bit ALU would be inefficient if not for vectorized dot instructions).

Another subtle co-design aspect is \textbf{memory format}. For instance, AMD’s GPUs prefer a tensor data format called NHWC (batch, height, width, channels) for some operations, while others prefer NCHW. Intel’s oneDNN library often uses blocked formats (like blocking channels in groups of 16 or 64) to better utilize AVX-512. These are software conventions chosen to match hardware capabilities (e.g., blocking by 16 for AVX-512 because 16 floats fit in a ZMM register). The software will transform data into these layouts to get faster computation, even though it’s extra work. Overall performance improves if the compute gain outweighs the transform cost.

In conclusion, software-hardware co-design in linear algebra is about making sure that the way we formulate computations in software (the algorithms, the data structures, the scheduling of operations) aligns well with what the hardware can do efficiently (the parallelism, the data movement, the precision, etc.). When done right, it can yield huge speedups. As evidence, deep learning models that once took days to train on CPUs can now train in hours on GPUs/TPUs, not only because the hardware is faster, but because every layer’s computation is carefully orchestrated by software to utilize that hardware (using batching, fusing operations, etc.). The constant feedback between software capabilities and hardware design (e.g., “We added feature X to chip because the compiler could then do Y more easily”) will continue to drive innovations in both fields.

\section{Future Directions in Processor Design for Linear Algebra}  
Looking forward, the demands of linear algebra-heavy workloads and the slowing of traditional CMOS scaling (Moore’s Law, Dennard scaling) are prompting exploration of novel computing paradigms. Here we outline a few speculative and emerging directions:

\subsection{Analog Computing and In-Memory Processing}  
One radical approach is to perform linear algebra operations using analog circuits rather than digital logic. The prime example is using analog electrical behavior to compute matrix-vector products directly in memory arrays. Consider a crossbar array of resistive memory (RRAM or MRAM): if you program the conductance of each crosspoint proportional to a matrix element, and input a vector as voltages on rows, then by Kirchhoff’s and Ohm’s laws the currents that accumulate on the columns represent the result of an analog matrix-vector multiplication (with sums happening as currents add up)【33†L35-L43】【33†L39-L47】. This is often called \textbf{analog in-memory computing} or a \textbf{memristor crossbar} approach. Research prototypes have shown that this can be “orders of magnitude more efficient than a highly optimized digital ASIC” for certain dot-product computations【33†L43-L47】, because you essentially get $N^2$ analog MAC operations happening in one step (all the crosspoints operate in parallel) and you only pay the energy to convert inputs to analog and read out outputs via ADCs.

The appeal is huge for linear algebra: one could potentially do a massive matrix multiplication in a single (analog) step, whereas digital would take many clock cycles. Also, it combines storage and compute, so data movement is minimized (the weights never leave the crossbar array). Analog computing is being pursued for neural network accelerators – effectively performing the multiply-accumulate of a layer with physical electronics. However, there are challenges: noise, limited precision (maybe 4-8 bits effectively), variability of devices, and analog circuits for addition are not as exact as digital. Nevertheless, for many ML applications some error is tolerable.

In terms of products, there are startups and prototypes of analog AI chips (using SRAM or RRAM technology for analog MACs). If these mature, we may see specialized analog co-processors for linear algebra that far outperform digital counterparts in energy efficiency. They might sit on the same chip as a CPU/GPU or be on separate accelerators.

Another approach is \textbf{optical computing}. Using light to compute has the advantage of potentially massive parallelism and high speed (speed of light and all). An optical matrix multiplier might use properties of light interference and diffraction to perform multiply-accumulate. For example, an optical system could encode a vector in a light wave and have it pass through a material (or network of waveguides) that modulates the light according to matrix weights; the resulting intensity pattern could represent the output vector. There have been demonstrations of optical dot-product engines and even optical neural networks. For instance, using Mach-Zehnder interferometers arrays to implement matrix multiplication for inference has been researched, and 3D optical computing (using volume holograms, etc.) promises extremely high throughput【48†L1-L4】.

Optics can naturally implement multiplication (modulating intensity or phase) and addition (combining beams). One advantage is that many frequencies (colors) of light can propagate simultaneously without interfering – \textbf{wavelength-division multiplexing} – which could be like doing many operations in parallel without additional hardware. Optical computing for AI is still in R\&D, but companies and labs are optimistic that it could overcome some limits of electronic computing, especially for matrix operations that demand both high bandwidth and low latency【37†L110-L118】【38†L7-L11】. A recent article noted that “optical computing performs calculations using the laws of physics… input information is encoded into properties of light, and through its propagation and interaction, the computation is achieved”【38†L7-L13】. The analog nature means results may need error correction or calibration.

\subsection{Quantum Accelerators}  
Quantum computing is fundamentally built on linear algebra – the state of a quantum system is a vector in a Hilbert space, and operations are matrices (unitary matrices for gates, for example). While general-purpose quantum computers aim to solve different problems (like factoring, Grover search, etc.), there is a subfield of \textbf{quantum algorithms for linear algebra}. Notably, the Harrow-Hassidim-Lloyd (HHL) algorithm provides a way to solve linear systems exponentially faster in principle than classical algorithms (giving solution vectors in quantum form)【25†L11-L18】. There are also quantum algorithms for tasks like Fourier transforms, finding eigenvalues, etc., which are linear algebra problems at their core.

In terms of hardware, quantum processors (qubit-based) are not going to execute a conventional matrix multiplication on classical data faster than a GPU anytime soon. However, they might be able to tackle certain large-scale linear algebra problems by working on encoded data. One could imagine a future where a quantum coprocessor is used for specific linear algebra subroutines (especially those where quantum algorithms offer polynomial or exponential speedups, like large linear system solving, least squares fitting, etc., for very large dimensions). But this is quite far out – current quantum devices are small and noisy.

Where quantum ideas might influence nearer-term hardware is in adopting quantum-inspired algorithms (like tensor network methods, which is a classical simulation technique for quantum systems that also doubles as a machine learning model approach for certain problems). Also, thinking of data in terms of quantum states has led to interesting analog processing ideas like coherent Ising machines for optimization (not directly linear algebra, but related to solving certain equations).

It’s safe to say that in the 10-20 year horizon, quantum accelerators will be a niche addition for specific tasks, and the bulk of linear algebra in everyday computing will still rely on classical hardware. But research continues; as one source puts it: many quantum algorithms “rely on linear algebra operations”【24†L31-L39】, so improvements in quantum hardware could one day radically alter how we approach big linear algebra problems.

\subsection{Neuromorphic and Brain-Inspired Computing}  
Neuromorphic chips (like Intel’s Loihi or IBM’s TrueNorth) are not directly linear algebra engines; they model spiking neural networks, which compute via event-driven spikes rather than dense linear algebra. However, one can approximate some linear algebra computations in a spiking paradigm (e.g., accumulations happening via spike rates). It’s a very different model and currently doesn’t outperform GPUs on standard linear algebra tasks. But if energy efficiency is a priority, brain-inspired event-based processing could handle certain vector-matrix operations extremely efficiently by only computing changes (sparse updates) rather than full updates every cycle. For instance, a neuromorphic system might naturally excel at multiplying a sparse matrix with a vector, if formulated as a network.

\subsection{3D Integration and Memory Proximity}  
One future trend is simply making current architectures more efficient by moving memory closer to compute. Technologies like HBM (High Bandwidth Memory) stacked on GPUs already improve things. Future CPUs/GPUs may use 3D stacking to put large chunks of SRAM (or even DRAM) directly above logic, enabling much higher bandwidth for feeding ALUs, which is very beneficial for linear algebra (since moving data is often the bottleneck). Also, there’s research into putting small compute units directly into memory chips (processing-in-memory, PIM). For example, a DRAM that can do a bulk bit-wise operation on a whole row of bits (this can be leveraged to do some linear algebra on bit-vectors very fast, or do matrix operations bit-serially but in huge parallel width).

In summary, the future likely holds a mix of:
- Continued evolution of digital architectures (wider vectors, more specialized matrix units, better memory systems, maybe more flexible “tensor cores” that handle multiple dimensions).
- Integration of analog/optical techniques for extremely fast and efficient matrix ops in domains where a bit of noise is acceptable (AI inference, maybe analog scientific computing solvers).
- Explorations of quantum for certain linear algebra problems.

From a computer architecture perspective, linear algebra will remain a driving application domain. We already see that in the design of new CPUs/GPUs (they highlight TOPS for AI, special accelerators, etc.). It’s reasonable to expect future general-purpose processors to look more like hybrids: part conventional, part matrix-engine, part(part conventional control logic, part specialized linear algebra engines, and perhaps even part analog/neuromorphic units for efficiency). Linear algebra will undoubtedly continue to be a driving force in computer architecture innovations, ensuring that future processors can meet the computational demands of scientific computing and AI alike.

\section{Conclusion}  

In this chapter, we explored the profound influence of linear algebra on computer processor design. We saw how the abstraction of vectors, matrices, and tensors is mirrored in hardware: from the vector registers and MAC units of DSPs, through the massively parallel matrix-crunching cores of GPUs, to the tensor-focused systolic arrays of TPUs. Historically, the quest to accelerate linear algebra computations led to vector supercomputers and SIMD extensions, and today the same quest is steering the industry toward specialized AI accelerators and novel computing paradigms. We delved into current CPU architectures (RISC-V, ARM, x86) to see concrete examples of linear algebra support — such as wide vector instructions, fused multiply-add operations, and even dedicated matrix instructions — and examined how operations like dot products and matrix multiplies are implemented efficiently on these platforms. We also discussed how higher-dimensional tensor operations can be decomposed or transformed to leverage those same efficient matrix/vector engines, highlighting a theme of "reduce to what hardware is good at."

A key takeaway is the importance of software/hardware co-design: performance gains in linear algebra often come from a tight coordination between algorithm, compiler, and processor capabilities. High-performance libraries and ML frameworks transparently harness low-level instruction sets and hardware features so that scientists and developers can get peak performance from high-level code. This synergy has enabled incredible progress — for example, enabling real-time computer vision and rapid neural network training that would be impractical without both sophisticated software optimization and powerful linear algebra hardware.

Finally, we cast an eye toward the future, considering how emerging technologies like analog in-memory computing, photonic processors, and quantum accelerators might further revolutionize linear algebra processing. While speculative, these approaches underscore a consistent trend: as long as vectors, matrices, and tensors remain central to computational problems, architects will keep inventing new ways to compute linear algebra operations faster and more efficiently. In understanding these processors — present and future — one essentially needs to understand linear algebra, as it is the language these machines speak at the most fundamental level.